\documentclass[10pt,a4paper]{amsart}
\usepackage[margin=19mm]{geometry}
\usepackage{amsmath,amssymb,mathtools}
\usepackage[scr=rsfs,cal=euler]{mathalfa}
\usepackage{xfrac}
\usepackage[unicode,hidelinks,bookmarks=false]{hyperref}
\newcommand{\ie}{i.e.\ }
\newcommand{\cf}{cf.\ }
\newcommand{\resp}{respectively\ }
\newcommand{\Subsection}{\subsection}
\providecommand{\nobreakdash}{\nobreak\hbox{-}}
\setlength{\emergencystretch}{2em}
\multlinegap0pt
\usepackage{bm}
\usepackage{bbm}
\usepackage{dsfont}
\usepackage{xspace}
\usepackage{cancel}
\usepackage{mathtools}
\usepackage{amsthm}
\makeatletter
\@ifundefined{thm}{\newtheorem{thm}{Theorem}}{}
\@ifundefined{lem}{\newtheorem{lem}[thm]{Lemma}}{}
\@ifundefined{prop}{\newtheorem{prop}[thm]{Proposition}}{}
\makeatother
\title[Frame bound, spectral gap and Plus space]{Frame bound, spectral gap and Plus space}
\author{Zheng-Yi Lu}
\date{Source: arXiv:2505.22136v1 (CC BY 4.0)}
\begin{document}
\maketitle

\noindent\emph{Source and licence.} Frame bound, spectral gap and Plus space. Version arXiv:2505.22136v1 (CC BY 4.0), distributed under the Creative Commons Attribution 4.0 International licence, as recorded in the arXiv licence metadata for this exact version and shown on the abstract page https://arxiv.org/abs/2505.22136v1. Full rights notice: licenses/arxiv-2505.22136-CC-BY-4.0.txt.\par\smallskip

\noindent\textit{Editorial scope.} Counted unit: the Prop whose source label is \texttt{thm4.5} in the article (source0.tex lines 690--693, source label \texttt{thm4.5}), delivered together with the article's own complete original proof of it (source0.tex lines 694--758, one \texttt{proof} environment of 65 lines). No mathematics is written, repaired or re-derived: statement and proof are the article's own text line for line. Required context retained uncounted: thm1.1 (lines 112--124); (1.2) (lines 206--208); thm4.1 (lines 545--548); thm4.2 (lines 613--617); thm4.3 (lines 634--636); thm4.4 (lines 654--660). Every definition, display and label of those lines is retained unchanged. Omitted from this package and not redistributed: 1--111, 125--205, 209--544, 549--612, 618--633, 637--653, 661--689, 759--911. Nothing inside a retained excerpt is omitted or altered: every retained body is delivered exactly as the article's lines are written, so no omission receipt of any kind accompanies this package. Citations: the retained excerpts carry no citation command at all, so no bibliography entry of the article is transcribed and no reference is redistributed. Reference adaptation: none was needed and none was made; every reference inside the delivered unit resolves inside it, so no unresolved reference or citation is produced. Printed slips kept literal: no printed word, symbol, hypothesis or number of the counted statement or of its proof was altered. Layout and class: the article's own class is \texttt{amsart} with packages including hyperref, amssymb, amsmath, amsthm, enumerate, graphicx, color, pgfplots, tikz; the wrapper is the lane's pinned \texttt{amsart} editorial wrapper at 10pt on A4, which restates the theorem-like environments the retained text needs and adds the article's own macro definitions that the retained text needs (none), with extra wrapper packages (bm, bbm, dsfont, xspace, cancel, mathtools). The delivered numbering is the unit's own. Assets: no retained span contains a figure environment, a graphics-inclusion command, a PDF-page inclusion, a table or a TikZ picture; no image file is copied into this package, the package declares no asset dependency and no image file is redistributed with it. Licence: version arXiv:2505.22136v1 of this article is distributed under the Creative Commons Attribution 4.0 International licence, as recorded in the arXiv licence metadata for this exact version and shown on the abstract page https://arxiv.org/abs/2505.22136v1. A full rights notice is delivered at licenses/arxiv-2505.22136-CC-BY-4.0.txt. Characters: every retained excerpt is pure ASCII, so the delivered engine, XeLaTeX with its default Latin Modern text font, reports no missing character.
\begin{thm}\label{thm1.1}
	Let $\mu$ be a Borel probability measure on $\Bbb R$. 
	Suppose that there exists $C>0$ such that \( |\widehat{\mu}(\xi)|\leq C|\xi|^{-1},\;\forall\xi\in\Bbb R\setminus\{0\}. \) 
If a countable set \(\Lambda\subset\Bbb R\) satisfies 
\begin{align}\label{(1.1')}
	\sum_{\lambda\in\Lambda}\Big|\int f(x)e^{-2\pi i\lambda x}d\mu(x)\Big|^2\geq A
	\int |f(x)|^2d\mu(x),\;\;\forall\;f(x)\in L^2(\mu),
\end{align}
then 
	\[
	g_{\min}(\Lambda) g_{\max}(\Lambda)\leq \frac{C^2\pi^2}{A}.
	\] 
\end{thm}

\begin{align}\label{(1.2)}
	\rho_{t_1,t_2}=\frac{1}{2}(\mathcal{L}_{[t_1,t_1+1]}\times\delta_0+\delta_0\times\mathcal{L}_{[t_2,t_2+1]}).
\end{align}

\begin{lem}\label{thm4.1}
Let $\rho_{t_1,t_2}$ be given as \eqref{(1.2)} and  $(\lambda_1,\lambda_2)\in\mathcal{Z}(\widehat{\rho_{t_1,t_2}})$. 
If $|\lambda_1|\leq0.8$, then $\lambda_1=\lambda_2$ or $\lambda_1+\lambda_2=0$. 
\end{lem}

\begin{lem}\label{thm4.2}
	Let $\rho_{t_1,t_2}$ be a spectral measure with a spectrum $\Lambda$. 
	Then $\tau_j(\Lambda)$ is a tight frame-spectrum of $\mathcal{L}_{[0,1]}$ with frame bounded $2$, where $$\tau_j(\Lambda):=\{\lambda_j:(\lambda_1,\lambda_2)\in\Lambda\}  
	$$ and $j=1,2$. 
\end{lem}

\begin{prop}\label{thm4.3}
	The Plus space $L^2(\rho_{-\frac12})$ admits no exponential orthogonal basis. 
\end{prop}

\begin{lem}\label{thm4.4}
		Let $\rho_{t_1,t_2}$ be a spectral measure with a spectrum $\Lambda$. 
		Suppose that $(t_1,t_2)\neq(-\frac12,-\frac12)$.
	Then there exists $K\in\Bbb Z\setminus\{0\}$ such that  $$\tau_1(\Lambda)=\tau_1(\Lambda)+K, 
	$$ 
	where $\tau_1(\Lambda)$ be given as Lemma  \ref{thm4.2}.
\end{lem}

\begin{prop}\label{thm4.5}
		Let $\rho_{t_1,t_2}$ be a spectral measure with a spectrum $\Lambda$. 
 Then $t_1-t_2\in\Bbb Z\setminus\{0\}$ or $t_1+t_2\in\Bbb Z\setminus\{-1\}$. 
\end{prop}

\begin{proof}
	From Proposition \ref{thm4.3}, we assume that $(t_1,t_2)\neq(-\frac12,-\frac12)$.
It follows from Lemma \ref{thm4.2} and Theorem \ref{thm1.1} that 
$g_{\min}(\tau_1(\Lambda))\leq \frac{\sqrt{2}}{2}$. 
By Lemma \ref{thm4.1} and the definition of $g_{\min}(\tau_1(\Lambda))$, with loss of generality, we assume $\mathbf{0},(\lambda_1,\lambda_2)\in\Lambda$ with $\lambda_1\in(0,0.8)$,  $\lambda_1\pm\lambda_2=0$ and 
\begin{align}\label{(4.3)}
	T(\lambda_1,\lambda_2)\in\Bbb Z\setminus2\Bbb Z.
\end{align} 
Lemma \ref{thm4.4} and $\mathbf{0}\in\Lambda$ allow us to assume that there exists $(n,m)\in\Lambda\setminus\{\mathbf{0}\}$. 
$$
	(-1)^{T(n-\lambda_1,m-\lambda_2)}\frac{\sin\pi(n-\lambda_1)}{n-\lambda_1}+	\frac{\sin\pi(m-\lambda_2)}{m-\lambda_2}=0.
$$
For each $(n,m)\in\Lambda\cap\Bbb Z^2$, by $(n-\lambda_1,m-\lambda_2)\in\mathcal{Z}(\widehat{\rho})$ and $|\lambda_1|=|\lambda_2|$, then 
\begin{align}\label{(4.4)}
|n-\lambda_1|=|m-\lambda_2|.
\end{align}
For any $(\lambda_1',\lambda_2')\in\Lambda\setminus\Bbb Z^2$, by $(\lambda_1',\lambda_2')-(n,m)\in\mathcal{Z}(\widehat{\rho_{t_1,t_2}})$, then 
\begin{align}\label{(4.5)}
\frac{\sin(\lambda_1'-n)\pi}{(\lambda_1'-n)\pi}\pm \frac{\sin(\lambda_2'-m)\pi}{(\lambda_2'-m)\pi}=0.
\end{align}
We prove the proposition by dividing two cases: $\lambda_1=\lambda_2$ and $\lambda_1=-\lambda_2$.

Suppose that $\lambda_1=\lambda_2$. 
If $\lambda_1=\frac12$, by \eqref{(4.3)}, one has 
$$
T(\lambda_1,\lambda_2)=t_1-t_2\in\Bbb Z\setminus2\Bbb Z.
$$
Now, we assume $\lambda_1\neq\frac12$, then \eqref{(4.4)} yields that $n=m$. 
Noting the choice of $(n,n)$ are infinitely many, denote as $\{(n_k,n_k)\}\subset\Lambda$ such that the symbols in \eqref{(4.5)} maintain uniform interpretation, then \eqref{(4.5)} becomes  
$$
\frac{\lambda_2'-n_k}{\lambda_1'-n_k}\sin\lambda_1'\pi\pm\sin\lambda_2'\pi=\lim_{k\rightarrow\infty}\frac{\lambda_2'-n_k}{\lambda_1'-n_k}\sin\lambda_1'\pi\pm\sin\lambda_2'\pi=\sin\lambda_1'\pi\pm\sin\lambda_2'\pi=0,
$$
which implies that $\lambda_1'=\lambda_2'$. 
This proves $\Lambda\subset\{(x,x):x\in\Bbb R\}$. 
This means $T(\lambda_1',\lambda_2')\in2\Bbb Z+1$ when $(\lambda_1',\lambda_2')\in\Lambda-\Lambda$ with $\lambda_1'\not\in\Bbb Z$. 
Then $\tau_1(\Lambda)\subset\Bbb Z\cup(\lambda_1+\Bbb Z)$. 
If $1\not\in\tau_1(\Lambda)$, by Lemma \ref{thm4.2}, then 
\begin{align*}
	2=	2\int_0^1|e^{-2\pi i x}|^2dx&=\sum_{\lambda\in(\Bbb Z\cup(\Bbb Z+\lambda_1))\cap\tau_1(\Lambda)}|\widehat{\mathcal{L}}_{[0,1]}(\lambda+1)|^2
	\\
	&=\sum_{\lambda\in(\Bbb Z+\lambda_1)\cap\tau_1(\Lambda)}|\widehat{\mathcal{L}}_{[0,1]}(\lambda+1)|^2\leq 1,
\end{align*}
a contradiction, as $\Bbb Z+\lambda_1$ is a spectrum of $\mathcal{L}_{[0,1]}$. 
So $1\in\tau_1(\Lambda)$, then 
$$
2(t_1-t_2)=T(1,1)=T(\lambda_1,\lambda_2)+T(1-\lambda_1,1-\lambda_2)\in2\Bbb Z,
$$
i.e., $t_1-t_2\in\Bbb Z$. 
By \eqref{(4.3)}, we know that $t_1\neq t_2$, which proves the case. 

Now, we consider the case $\lambda_1=-\lambda_2$. 
If $\lambda_1=\frac12$, by \eqref{(4.3)}, one has 
$$
T(\lambda_1,\lambda_2)=t_1+t_2+1
\in\Bbb Z\setminus2\Bbb Z.
$$
Now, we assume $\lambda_1\neq\frac12$, then \eqref{(4.4)} yields that $n=-m$. 
Similar to the arguments of case $\lambda_1=\lambda_2$, we can get $\Lambda\subset\{(x,-x):x\in\Bbb R\}$,  $\tau_1(\Lambda)\subset\Bbb Z\cup(\lambda_1+\Bbb Z)$ and  $1\in\tau_1(\Lambda)$. 
Hence
$$
2(t_1+t_2+1)=T(1,1)=T(\lambda_1,\lambda_2)+T(1-\lambda_1,1-\lambda_2)\in2\Bbb Z,
$$
i.e., $t_1+t_2+1\in\Bbb Z$. 
Finally, using \eqref{(4.3)} again, we reduce  that $t_1+t_2\neq-1$, which proves the proposition. 
\end{proof}

\end{document}
